% TOP_PROBLEM: {"id": 3383, "title": "Divisibility of central binomial coefficients", "category_id": 1, "category": {"id": 1, "name": "number_theory", "display_name": "Number Theory", "description": "Properties of integers, prime numbers, Diophantine equations.", "slug": "number-theory", "order_index": 1, "created_at": "2026-07-31T15:26:25.670Z"}, "status": "open"}
% FIRST_POSTED: null
% ENTRY_KIND: statement_only
% Imported from: batch04.tex; UnsolvedMath ID 3383
\documentclass[11pt]{article}
\usepackage[margin=25mm]{geometry}
\usepackage{fontspec}
\setmainfont{Latin Modern Roman}
\usepackage{amsmath,amssymb,mathtools}
\usepackage{unicode-math}
\setmathfont{Latin Modern Math}
\usepackage{xurl,hyperref}
\hypersetup{colorlinks=true,urlcolor=blue}
\setcounter{secnumdepth}{0}
\setlength{\parindent}{0pt}
\setlength{\parskip}{5pt}
\setlength{\emergencystretch}{3em}
\newcommand{\sref}[2]{\hyperlink{source:#1:#2}{[S#2]}}
\newcommand{\cataloguescope}[1]{\paragraph{Recorded target.}#1}
\begin{document}
% UnsolvedMath ID 3383; source catalogue code OPG-506
\setcounter{section}{3382}
\section{Divisibility of central binomial coefficients}\label{3383}
{\small\sffamily UnsolvedMath ID 3383 · Number Theory}
\subsection{Definitions and mathematical statement}

\paragraph{Shared notation.}$\mathbb N=\{1,2,\ldots\}$, and $\mathbb Z,\mathbb Q,\mathbb R,\mathbb C$ denote the integers, rationals, reals and complex numbers. Any other conventions are specified locally.


For $n\ge1$, let $B_n=\binom{2n}{n}=(2n)!/(n!)^2$ and let $\gcd$ denote the greatest common divisor of positive integers.
\paragraph{Statement recorded in the supplied source.}
Settle both assertions:
\begin{enumerate}
\item The set $\{n\ge1:\gcd(B_n,3\cdot5\cdot7)=1\}$ is infinite.
\item The set $\{n\ge1:\gcd(B_n,3\cdot5\cdot7\cdot11)=1\}$ is finite.
\end{enumerate}
For an odd prime $p$, the condition $p\nmid B_n$ is equivalently that every digit of $n$ in base $p$ is at most $(p-1)/2$. \sref{3383}{2}
\subsection{Short English statement}
Determine whether avoidance of the primes three, five and seven occurs infinitely often for central binomial coefficients, while simultaneous avoidance of eleven occurs only finitely often.
\subsection{Sources}
\begin{description}
\item[\hypertarget{source:3383:1}{S1}] ULAM.AI and the UnsolvedMath Contributors, UnsolvedMath: A Curated Collection of Open Mathematics Problems, record 3383 (OPG-506). CC BY 4.0. \url{https://huggingface.co/datasets/ulamai/UnsolvedMath}.
\item[\hypertarget{source:3383:2}{S2}] Open Problem Garden, Divisibility of central binomial coefficients, OPG-506. Original problem statement, full mathematical discussion, bibliography and comments. \url{https://www.openproblemgarden.org/op/divisibility_of_central_binomial_coefficients}.
\end{description}
\label{end:3383}
\end{document}
